\vssub
\subsection{~波浪模式例程} \label{sec:core}
\vssub

波浪模式驱动器是 \ws{} 框架包中的一个子程序。要运行该模式驱动子程序，需要一个程序壳。\ws{} 提供了一个简单的独立壳，见 \para\ref{sec:ww3shel}，也提供了一个更复杂的多网格模式壳，见 \para\ref{sec:ww3multi}。本节重点介绍波浪模式驱动子程序。

波浪模式初始化例程 {\F w3init} 对单个波浪模式网格执行模式初始化。这包括：通过定义单元号来建立部分 I/O 系统；初始化内部时间管理；处理模式定义文件（{\file mod\_def.ww3}）；处理初始条件（{\file restart.ww3}）；准备模式输出；以及计算依赖网格的参数。若模式为 \mpi{} 环境编译，则会确定并初始化计算和输出所需的全部通信（模式全程使用持久 \mpi{} 通信）。

在初始化完成后，可任意多次调用波浪模式例程 {\F w3wave}，以推进单个网格上的波场随时间演变。在若干初始检查之后，该子程序会插值风和流、更新冰浓度和水位、传播波场，并在若干时间步上应用所选源项。内部时间步由需要执行计算的时间区间以及请求的输出时刻共同定义。计算结束时，该例程向调用程序提供所请求的波浪数据场。{\F w3wave} 接口的文档可在源代码（{\file w3wavemd.ftn}）中找到。

\input{zh/run/fig_log_zh.tex} 

除上述原始数据文件外，程序还维护一个日志文件 {\file log.ww3}。该文件由 {\F w3init} 打开（{\F w3init} 包含在 {\file w3wavemd.ftn} 的 {\F w3wave} 中），并向该文件写入若干自解释的头部信息。对 {\F w3wave} 的每次后续调用都会在该日志文件的“动作表”中增加数行，如图~\ref{fig:log} 所示。标识为 `step' 的列显示所考虑的离散时间步。标识为 `pass' 的列给出调用 {\F w3wave} 的序列号；例如 3 表示该动作发生在第三次调用 {\F w3wave} 中。第三列显示该时间步的结束时间。输入和输出列显示模式的相应动作。{\tt X} 表示输入已更新或输出已执行。{\tt F} 表示首次读取场，{\tt L} 表示最后一次输出。七个输入列分别表示边界条件（{\tt b}）、风场（{\tt w}）、水位（{\tt l}）、流场（{\tt c}）、冰浓度（{\tt i}）以及同化数据（{\tt d}）。注意，数据同化发生在波浪例程调用之后、时间步末端。七个输出列分别表示网格输出（{\tt g}）、点输出（{\tt p}）、沿轨迹输出（{\tt t}）、重启动文件（{\tt r}）、边界数据（{\tt b}）、分区谱数据（{\tt f}）以及耦合输出（{\tt c}）。

对于多网格波浪模式 \citep[][{\file ww3\_multi}]{tol:OMOD08b}，在基本波浪模式例程周围构建了一组例程。三个主要例程是初始化例程 {\F wminit}、时间步进例程 {\F wmwave} 和终止例程 {\F wmfinl}，其功能与上述单网格例程类似。注意，原始输入和输出文件会针对镶嵌中的各个网格分别生成，并通过将标准文件扩展名 `{\file .ww3}' 替换为用户在 {\file ww3\_grid.inp} 文件中为每个单独网格选择的唯一标识符来识别。每个单独网格都维护日志文件，同时也维护一个总日志文件 {\file log.mww3}。
